\section{Omniprediction for multiple quantile level forecasts}\label{sec:multi_q}
In this section, we extend the result from Section \ref{sec:single_q} to multiple quantile levels. % To avoid redundancy, we mainly provide online algorithm for omniprediction of multiple quantile level forecasts. Batch version can be straightforwardly extended from the single quantile level case. 
The omniprediction objective we want to minimize is given by
\[
\sup_{S \in \cs,\, f \in \cf} \frac{1}{T} \sum_{t=1}^T \E_{p \sim \hat{P}} [S(p, Y_t)] - \E[S(f_t, Y_t)]
\]
where $\cs$ is the class of \textit{all} scoring functions that are consistent for the multiple quantile levels and $\cf$ is a finite set of competitors or base forecasters.

We follow a similar flow as in Section \ref{sec:single_q}. Section \ref{subsec:decomp_multi_q} discusses the previous result from \cite{fissler2016multiq} on decomposition of scoring function that is consistent for multiple quantile levels. With this result and same discretization scheme from Section \ref{subsec:discretization}, Section \ref{subsec:omni_multi_q} derive a closed-form solution to the multiple quantile level version of minimax program \eqref{eq:minimax_Pt_multi_q} and present the multiple quantile level omniprediction algorithm. Lastly, Section \ref{subsec:separate_multi_q} compares multiple quantile level omniprediction with running separate single quantile level omniprediction for each quantile level.

%Section \ref{subsec:separate_multi_q} shows the cases where post-sorting after solving separate minimax programs for each quantile level works and the cases where it fails, which necessitates the need for solving the minimax program \eqref{eq:omni-twoplayer-minimax} jointly across all quantile levels. 
% Section \ref{subsec:omni_multi_q} provides a closed form solution to the minimax program \eqref{eq:omni-twoplayer-minimax} for multiple quantile levels. Finally, Section \ref{subsec:omni_multi_q_online} puts the results together and provides a sample efficient algorithm for omniprediction of multiple quantile level forecasts.

% \cite{fissler2016multiq} proved that a consistent (proper) loss for multiple quantile levels can be written as a weighted sum of consistent (proper) loss for each quantile levels, under mild regularity conditions.


\subsection{Decomposition of score functions that are consistent for the multiple quantile levels}\label{subsec:decomp_multi_q}
In this subsection, we review the structural characterization of scoring functions that are consistent for multiple quantile levels. \cite{fissler2016multiq} shows that, under mild regularity conditions, any scoring function that is consistent for multiple quantile levels can be written as a weighted sum of scoring functions that are consistent for each individual quantile level. For further details, see Proposition 4.1 and Corollary 4.2 in \cite{fissler2016multiq}. Here, we provide a simplified version of these results that covers only the scope needed for this paper.

\begin{theorem}[Simplified version of Corollary 4.2 in \cite{fissler2016multiq}]
Let $\cf_0$ be the class of the probability measures on the Borel-Lebesgue sets of the real line $\R$. Suppose that $T = (q_{\tau_1}, \ldots, q_{\tau_N}): \cf_0 \to \R^N$ is a functional where each $q_{\tau}$ is a quantile functional at level $\tau$. Then, up to mild regularity conditions, a scoring function $S: \R^N \times \R \to \R$ is strictly consistent for $T$ relative to $\cf_0$ if and only if it is of the form
\begin{align}
    S(x_1, \ldots, x_N, y) = \sum_{n=1}^{N} S_n(x_n, y),
\end{align}
for almost all $(\{x_n\}_{n=1}^N, y) \in \R^N \times \R$, where $S_n: \R \times \R \to \R$, $n \in [N]$, are strictly consistent scoring functions for $q_{\tau_n}$ relative to $\cf_0$.
\end{theorem}

Combined with the decomposition of scoring functions that are consistent for a single quantile level (Theorem \ref{thm:decomp_single_q}), a scoring function that is consistent for multiple quantile levels, $\{\tau_n\}_{n=1}^N$, can be written as a weighted sum of elementary scores of each quantile level. 
\begin{align}\label{eq:multi_q_decomp_raw}
S(x_1, \ldots, x_N, y) = \sum_{n=1}^{N} \int_{-\infty}^{\infty} S_{\tau_n, \theta}(x_n, y) dH_{n}(\theta),
\end{align}
where $H_n$ is a nonnegative measure and $S_{\tau_n, \theta}(x_n, y) = (\one(y<x_n)-\tau_n)(\one(\theta<x_n)-\one(\theta<y))$ for $n \in [N]$ and $\theta \in \R$. 





\subsection{Omniprediction algorithm for multiple quantile levels} \label{subsec:omni_multi_q}
We begin this subsection by introducing the notations and then present the algorithm for multiple quantile level omniprediction. We derive the closed-form solution to the minimax program \eqref{eq:minimax_Pt_multi_q} and present the theorem that provides an upper bound for the omniprediction error of the algorithm.

\paragraph{Notations}
We consider quantile levels $\{\tau_n\}_{n=1}^N$ where $\tau_1 < \tau_2 < ... < \tau_N$ and each base forecasters $f_b$, $b \in [B]$ submits an $N$-length forecast $\{f_b^n\}_{n=1}^N$ for each time step $t=1, ..., T$, where $f_b^n(t)$ corresponds to the $\tau_n$-level quantile estimate from $f_b$ at time $t$.

\paragraph{Discretization}
In a similar vein to the discretization scheme introduced in Section \ref{subsec:discretization}, without loss of generality, we consider the ground truth $Y_t$ and base forecasters $\{f_b^n(t)\}$, $b \in [B], n \in [N]$ to take values in the set $\{0, \frac{1}{m}, ..., 1\}$ and use the same grid of $\{\theta_i\}_{i=1}^{m}$, where $\theta_i = \frac{1}{2m} - \frac{i}{m}$ for $i \in [m]$. Combined with the characterization \eqref{eq:multi_q_decomp_raw}, we can approximate a scoring function that is consistent for multiple quantile levels, $S$, as a weighted sum of elementary scores of each quantile level.
\begin{align*} %\label{eq:multi_q_decomp_discretized}
  S(x_1, \ldots, x_N, y) = \sum_{n=1}^{N} \sum_{i=0}^{m-1} w_i^n \cdot S_{\tau_n, \theta_i}(x_n, y),
\end{align*}
where $w_i^n$ is the weight for the $i$-th grid point of the $n$-th quantile level. 

With the discretization, Algorithm \ref{alg:omni-twoplayer_online_multi_q} presents the multiple quantile level omniprediction algorithm. We abbreviate $S_{\tau_n, \theta_i}$ as $S_{n, i}$.

\begin{algorithm}
  \caption{Two-player game based multiple quantile level online omniprediction} \label{alg:omni-twoplayer_online_multi_q}
  \begin{algorithmic}[1]
  \State \textbf{Data:} Data stream $\{Y_t\}_{t=1}^T$, quantile levels $\tau_1 < \tau_2 < ... < \tau_N$, Grid size $m\in \mathbb{N}$, Learning rate $\eta> 0$, Base forecasters: $\{f_b\}_{b=1}^B$ that outputs $f_b^n(t)$ as $\tau_n$-level quantile estimate at time $t$;
  \State \textit{Initialize:} $w_{ib}^n(1) = \frac{1}{mNB}$, $\forall i \in [m], n \in [N], b \in [B]$;
  \For{$t = 1, \ldots, T$}
    \State \textbf{1) Make prediction at time $t$}:
      \State Define $V_0 = V_N = 0$, $j_0^* = 0$, $j_N^* = m$;
      \State Define $\one_i^n(t) = \{\indi{\theta_i < f_b^n(t)}\}_{b=1}^B \in \R^B$ and $\one_B = (1,..,1)\in \R^B$ for $i \in [m]$ and $n \in [N]$.
      \For{$n=1, \ldots, N-1$}
      \State Compute $V_n(j)$ for $j \in [m]$:
        \(
        V_n(j) = \sum_{i=j+1}^{m} \sum_{s=1}^{n} w_i^s \cdot \one_i^s(t) + \sum_{i=1}^{j} \sum_{s=n+1}^{N}  w_i^s \cdot (\one_B - \one_i^s(t))
        \)
      \State Compute $V_n$ and $j_n^*$:
        \begin{align*}
          V_n = \min_{j \in [m]} V_n(j), \qquad
          j_n^* = \min\left\{j \mid j\in [j_{n-1}^*, m],  V_n(j) = V_n \right\}
        \end{align*}
      \EndFor \vspace{-0.2cm}\\
      \For{$n=1, \ldots, N$}
        \If{$V_n(j_{n-1}^*) = V_n(j_n^*)$}
          \State $\hat{P}_t^n = \delta_{j_{n-1}^*/m}$;
        \Else
          \State Define $B_k := \sum_{i=1}^{k} w_i^n(t) \cdot \one_B - \sum_{i=1}^{m}w_i^n(t) \cdot \one_i^n(t)$
          \State Find $j_{n-1}^* \leq k_n < j_n^*$ s.t. $B_{k_n} < V_{n-1} - V_n \leq B_{k_n+1}$ (Existence proved in Lemma \ref{lem:Pn_def})
          \State $\hat{P}_t^n = (1-C_n)\delta_{k_n/m} + C_n\delta_{((k_n+1)/m)}$, where $C_n = \left( V_{n-1} - V_n - B_{k_n} \right) / (w_{k_n}^n \cdot \one_B)$
        \EndIf
      \EndFor
      \State \textbf{Output} $\hat{P}_t = \{\hat{P}_t^n\}_{n=1}^N$;  \vspace{-0.1cm}\\
      \State \textbf{2) Update the weights for each $S_{n,i} = (\one(y<x)-\tau_n)(\one(\theta_i<x)-\one(\theta_i<y))$}:
          \State $w_{ib}^n(t+1) 
          = \frac{w_{ib}^n(t) \exp\bigl(\eta \big(\mathbb{E}_{p \sim \hat{P}_t^n}[S_{n,i}(p, Y_t)] 
          - S_{n,i}(\hat{f}_i^n(t), Y_t)\big)\bigr)}
          {\sum_{i=1}^{m} \sum_{n=1}^N \sum_{b=1}^B w_{ib}^n(t) \exp\bigl(\eta \big(\mathbb{E}_{p \sim \hat{P}_t^n}[S_{n,i}(p, Y_t)] - S_{n,i}(\hat{f}_b^n(t), Y_t)\big)\bigr)}, \quad \forall i \in [m], n \in [N], b \in [B];$ \vspace{-0.1cm}\\
  \EndFor
  \end{algorithmic}
\end{algorithm}



%%% HERE Apr. 13

As in the case of single quantile level omniprediction, the first part of the algorithm is to make a prediction at time $t$. $\hat{P}_t$ computed here is a minimax solution for the following minimax program, which hasone additional dimension compared to the single quantile level case:

\[
\min_{\tilde{P} \in \cd(\Delta_m)} \max_{P_y \in \Delta_m} \sum_{n=1}^N \sum_{i=0}^{m-1} \sum_{b=1}^F w_{ib}^n \E_{P\sim \tilde{P}, Y \sim P_y}[S_i^n(q_{\tau_n}(P),Y) - S_i^n(f_b^n(t),Y)],
\]
where $q_{\tau_n}(P)$ is the $\tau_n$-level quantile of $P$ and $S_i^n(x, Y) = (\one(Y < x)-\tau_n)(\one(\theta_i<x)-\one(\theta_i<Y))$ as defined previously. We can rewrite the program as follows:
\begin{align}\label{eq:minimax_Pt_multi_q}
\min_{P_1, P_2, ..., P_N \in \Delta_m} \max_{P_y \in \Delta_m} \sum_{n=1}^N \sum_{i=0}^{m-1} \sum_{b=1}^F w_{ib}^n \E_{p_n \sim P_n
, Y \sim P_y}[S_i^n(p_n,Y) - S_i^n(f_b^n(t),Y)],
\end{align}
with the constraint that the joint distribution ($P_1, P_2, ..., P_N$) always satisfies the ordering of the samples $p_1 \leq p_2 \leq ... \leq p_N$. During derivation of the closed-form solution, we also prove that our closed-form solution satisfies this monotonicity constraint.


\begin{lemma}\label{lem:jn_sort}
Define the quantity $V_n$ for $n \in [N-1]$ as
\begin{align} \label{eq:Vn_def}
V_n := \min_{j \in [m]} \left( \sum_{i=j}^{m-1} \sum_{s=1}^{n} \sum_{b=1}^F w_{ib}^s \indi{\theta_i < f_b^s(t)} + \sum_{i=0}^{j-1} \sum_{s=n+1}^{N}  \sum_{b=1}^F w_{ib}^s (1 - \indi{\theta_i < f_b^s(t)})\right) \geq 0
\end{align}
and $V_0 = V_N := 0$.
Let $J_n^*$ be the set of indices $j \in [m]$ that minimizes the $V_n$. 
\smallskip \\
Then, there exists a sorted sequence $0=j_0^* \leq j_1^* \leq ... \leq j_{N-1}^* \leq j_N^* = m$ such that $j_n^* \in J_n^*$ for $n \in [N-1]$.
\end{lemma}

\begin{proof}
Define $w_i^s := (w_{1b}^s, ..., w_{Fb}^s) \in \R^F$, $\one_i^s = (\indi{\theta_i < f_1^s(t)}, ..., \indi{\theta_i < f_F^s(t)}) \in \R^F$, and $\one_F = (1, ..., 1) \in \R^F$. Then, $V_n$ can be rewritten as
\[
V_n = \min_{j \in [m]} \left( \sum_{i=j}^{m-1} \sum_{s=1}^{n} w_i^s \cdot \one_i^s + \sum_{i=0}^{j-1} \sum_{s=n+1}^{N}  w_i^s \cdot (\one_F - \one_i^s) \right).
\]
Suppose there exists a pair of indices $n, n+1$ such that $j_n^* > j_{n+1}^*$. Using the definition of $V_n$, we have 
\begin{align*}
j_n^* &= \argmin_{j \in [m]} \left( \sum_{i=j}^{m-1} \sum_{s=1}^{n} w_i^s \cdot \one_i^s + \sum_{i=0}^{j-1} \sum_{s=n+2}^{N}  w_i^s \cdot (\one_F - \one_i^s) + \sum_{i=0}^{j-1} w_i^{n+1} \cdot (\one_F - \one_i^{n+1})\right) \\
j_{n+1}^* &= \argmin_{j \in [m]} \left( \sum_{i=j}^{m-1} \sum_{s=1}^{n+1} w_i^s \cdot \one_i^s + \sum_{i=0}^{j-1} \sum_{s=n+2}^{N}  w_i^s \cdot (\one_F - \one_i^s) + \sum_{i=j}^{m-1} w_i^{n+1} \cdot \one_i^{n+1}\right).
\end{align*}
Two expressions share the first two terms and differ in the last term. We abbreviate these terms as following:
\begin{align*}
C_j & := \sum_{i=j}^{m-1} \sum_{s=1}^{n} w_i^s \cdot \one_i^s + \sum_{i=0}^{j-1} \sum_{s=n+2}^{N}  w_i^s \cdot (\one_F - \one_i^s) \\
D_j & := \sum_{i=0}^{j-1} w_i^{n+1} \cdot (\one_F - \one_i^{n+1}) \\
E_j & := \sum_{i=j}^{m-1} w_i^{n+1} \cdot \one_i^{n+1}.
\end{align*}
Note that $D_j$ is non-decreasing in $j$ and $E_j$ is non-increasing in $j$.
Since $j_n^*$ and $j_{n+1}^*$ are the indices that minimize the $V_n$ and $V_{n+1}$, respectively, we have $C_{j_n^*} + D_{j_n^*} \leq C_{j_{n+1}^*} + D_{j_{n+1}^*}$ and $C_{j_n^*} + E_{j_n^*} \geq C_{j_{n+1}^*} + E_{j_{n+1}^*}$. This implies that 
\[
E_{j_{n+1}^*} - E_{j_n^*} \leq C_{j_n^*} - C_{j_{n+1}^*} \leq D_{j_{n+1}^*} - D_{j_n^*}.
\]
Because of the monotonicity of $D_j$ and $E_j$ and the assumption that $j_n^* > j_{n+1}^*$, the leftmost term should be non-negative and the rightmost term should be non-positive. Therefore, we obtain 
\[
C_{j_n^*} - C_{j_{n+1}^*} = D_{j_n^*} - D_{j_{n+1}^*} = E_{j_n^*} - E_{j_{n+1}^*} = 0 \quad \Rightarrow j_{n+1}^* \in J_n^*, \,\, j_n^* \in J_{n+1}^*.
\]
By choosing $j_{n+1}^*$ as a minimizer of $V_n$ and $j_n^*$ as a minimizer of $V_{n+1}$, we resolve the unsorted pair of indices. Repeating this process yields a sorted sequence $0=j_0^* \leq j_1^* \leq ... \leq j_{N-1}^* \leq j_N^* = m$ and this proves the existence of sorted sequence.
\\
Due to its existence, we can simply obtain a sorted sequence by choosing $j_n^* = \min J_n^*$ for all $n \in [N]$.
\end{proof}

\begin{lemma} \label{lem:Pn_def}
Given a sorted sequence $0=j_0^* \leq j_1^* \leq ... \leq j_{N-1}^* \leq j_N^* = m$, there exists a distribution $P_n^*$ for $n \in [N]$ of the form
\begin{align}\label{eq:Pn_def}
P_n^* = (1-C_n)\delta_{k_n / m} + C_n\delta_{(k_n+1)/m},
\end{align}
such that 
\begin{enumerate}
  \item $j_{n-1}^* \leq k_n$
  \item $k_n < j_n^*$ if $j_{n-1}^* \neq j_n^*$ and $k_n = j_n^*$ with $C_n=0$ if $j_{n-1}^* = j_n^*$.
  \item The support of $P_n^*$ is contained in $[j_{n-1}^*/m, j_n^*/m]$.
  \item $\E_{p_n \sim P_n^*} \left[ \sum_{i=0}^{m-1} \sum_{b=1}^F w_{ib}^n \left( \indi{\theta_i < p_n} - \indi{\theta_i < f_b^n(t)} \right) \right] = V_{n-1}-V_n$.
\end{enumerate}
\end{lemma}
\begin{proof}
We again use the notations used in proof of Lemma \ref{lem:jn_sort}. We define a function $V_n(\cdot)$ as following:
\begin{align}\label{eq:Vn_j_def}
V_n(j) = \sum_{i=j}^{m-1} \sum_{s=1}^{n} w_i^s \cdot \one_i^s + \sum_{i=0}^{j-1} \sum_{s=n+1}^{N}  w_i^s \cdot (\one_F - \one_i^s)  \quad \forall j \in [m], n \in [N].
\end{align}
Note that the quantity $V_n$ defined in Lemma \ref{lem:jn_sort} is equal to $V_n(j_n^*)$ for all $n \in [N-1]$ and that $V_n(j) \geq V_n(j_n^*)$ for all $j \in [m]$.
\\
With the fact that $\indi{\theta_i < k/m} = \indi{i < k}$, the last condition in the lemma statement is equivalent to 
\begin{align*} 
& \E_{p_n \sim P_n^*} \left[ \sum_{i=0}^{m-1} \sum_{b=1}^F w_{ib}^n \left( \indi{\theta_i < p_n} - \indi{\theta_i < f_b^n(t)} \right) \right] \nonumber\\
& = (1-C_n) \sum_{i=0}^{m-1} \sum_{b=1}^F w_{ib}^n \left( \indi{i < k_n} - \indi{\theta_i < f_b^n(t)} \right) + C_n \sum_{i=0}^{m-1} \sum_{b=1}^F w_{ib}^n \left( \indi{i < k_n+1} - \indi{\theta_i < f_b^n(t)} \right) \nonumber\\
& = (1-C_n) \sum_{i=0}^{m-1} w_{i}^n \cdot \left( \indi{i < k_n} \one_F - \one_i^n \right) + C_n \sum_{i=0}^{m-1} w_{i}^n \cdot \left( \indi{i < k_n+1} \one_F - \one_i^n \right) \nonumber\\
& = \sum_{i=0}^{k_n-1} w_i^n \cdot \one_F + C_n w_{k_n}^n \cdot \one_F - \sum_{i=0}^{m-1} w_i^n \cdot \one_i^n \\
& = V_{n-1} - V_n.  
\end{align*}
The last equality is equivalent to 
\begin{align}
\sum_{i=0}^{k_n-1} w_i^n \cdot \one_F + C_n w_{k_n}^n \cdot \one_F 
= \sum_{i=0}^{m-1} w_i^n \cdot \one_i^n + V_{n-1} - V_n 
= \sum_{i=0}^{j_{n-1}^*-1} w_i^n \cdot \one_F + V_n(j_{n-1}^*) - V_n(j_n^*), \label{eq:Pn_eq}
\end{align}
where we used the fact that $\sum_{i=0}^{m-1} w_i^n \cdot \one_i^n + V_{n-1}(j) = \sum_{i=0}^{j-1} w_i^n \cdot \one_F + V_n(j)$ with $j=j_{n-1}^*$.
\smallskip \\
Now we divide the cases into three cases. 
First, consider the case where $j_{n-1}^* = j_n^*$. In this case, simply $k_n = j_{n-1}^*$ with $C_n=0$ satisfies the equation \eqref{eq:Pn_eq}. Therefore, $P_n^* = \delta_{j_{n-1}^*/m}$ satisfies the conditions.
Second, consider the case where $j_{n-1}^* < j_n^*$ but $V_n(j_{n-1}^*) = V_n(j_n^*)$. $k_n = j_{n-1}^*$ with $C_n=0$ again satisfies the equation \eqref{eq:Pn_eq} and $k_n = j_{n-1}^* < j_n^*$ in this case.
Lastly, consider the case where $V_n(j_{n-1}^*) > V_n(j_n^*)$. We have 
\begin{align*}
  \sum_{i=0}^{j_{n-1}^*-1} w_i^n \cdot \one_F
  & < \sum_{i=0}^{j_{n-1}^*-1} w_i^n \cdot \one_F + V_n(j_{n-1}^*) - V_n(j_n^*) \\
  & = \sum_{i=0}^{m-1} w_i^n \cdot \one_i^n - V_n + V_{n-1} \\
  & = \sum_{i=0}^{j_{n}^*-1} w_i^n \cdot \one_F + V_{n-1}(j_{n-1}^*) - V_{n-1}(j_n^*) 
  \leq \sum_{i=0}^{j_{n}^*-1} w_i^n \cdot \one_F,
\end{align*}

where two equalities used the fact that $\sum_{i=0}^{m-1} w_i^n \cdot \one_i^n + V_{n-1}(j) = \sum_{i=0}^{j-1} w_i^n \cdot \one_F + V_n(j)$ with $j=j_n^*$ on the second term of the equation \eqref{eq:Pn_eq}.
This implies that there exists a $j_{n-1}^* \leq k_n < j_n^*$ such that 
\begin{align*}
  \sum_{i=0}^{k_n} w_i^n \cdot \one_F \geq \sum_{i=0}^{m-1} w_i^n \cdot \one_i^n - V_n + V_{n-1} \quad \text{and} \quad \sum_{i=0}^{k_n-1} w_i^n \cdot \one_F < \sum_{i=0}^{m-1} w_i^n \cdot \one_i^n - V_n + V_{n-1}.
\end{align*}
With such $k_n$, we define $P_n^*$ with $C_n = \frac{1}{w_{k_n}^n \cdot \one_F} \left( \sum_{i=0}^{m-1} w_i^n \cdot \one_i^n - V_n + V_{n-1} - \sum_{i=0}^{k_n-1} w_i^n \cdot \one_F \right)$ and this satisfies the conditions since $k_n+1 \leq j_n^*$.
\end{proof}





Now we provide the solution to the minimax program \eqref{eq:omni-twoplayer-minimax}.
\begin{theorem}\label{thm:minimax_negative_multi_q}
Given quantile levels $\tau_1 < \tau_2 < ... < \tau_N$, quantile forecasts $\hat{f}_i^n$ for $i =0,...,m\!-\!1$ and $n \in [N]$, define $V_n$ $n \in [N-1]$ as in equations \eqref{eq:Vn_def}, with $V_0 = V_N := 0$. Let $\{j_n^*\}_{n=1}^{N-1}$ be a sorted sequence of indices that minimizes the $V_n$ for $n \in [N-1]$ with $j_0^* = 0$ and $j_N^* = m$ where the existence of the sequence is guaranteed by Lemma \ref{lem:jn_sort}. 
\smallskip \\
Then, $P_n^*$ for $n \in [N]$ given by \eqref{eq:Pn_def} with 
\[
P_y^* = \tau_1 \delta_0 + (1-\tau_N)\delta_1 + \sum_{n=1}^{N-1}(\tau_{n+1}-\tau_n)\delta_{j_t^*/m}
\]
solves the minimax program \eqref{eq:omni-twoplayer-minimax} and the value of the minimax program is 
\[
-\sum_{n=1}^{N-1}(\tau_{n+1}-\tau_n)V_n \leq 0.
\]

% Then, the value of the minimax program is $-\sum_{t=1}^{N-1}(\tau_{t+1}-\tau_t)V_t \leq 0$ and the solution is given by 
% \begin{align*}
%   P_y^* &= \tau_1 \delta_0 + (1-\tau_N)\delta_1 + \sum_{t=1}^{N-1}(\tau_{t+1}-\tau_t)\delta_{j_t^*/m} \\
%   P_t^* &= C_t \delta_{\underline{k_t^*}/m} + (1-C_t) \delta_{(\underline{k_t^*}+1)/m}, \quad t \in [N],
% \end{align*}
% where
% \begin{align*}
%   \underline{k_t^*} &= \min\left\{k \mid \sum_{i=0}^k w_i^t \geq \sum_{i=0}^{m-1} w_i^t \indi{\theta_i < \hat{f}_i^t} - V_t + V_{t-1} \right\} \\
%   \overline{k_t^*} &= \min\left\{k \mid \sum_{i=0}^k w_i^t > \sum_{i=0}^{m-1} w_i^t \indi{\theta_i < \hat{f}_i^t} - V_t + V_{t-1} \right\} \\
%   C_t & = \frac{1}{w_{\underline{k_t^*}}^t} \left( \sum_{i=0}^{\underline{k_t^*}} w_i^t - \left(\sum_{i=0}^{m-1} w_i^t \indi{\theta_i < \hat{f}_i^t} - V_t + V_{t-1}\right)\right),
% \end{align*}
% for $t \in [N]$ with $V_0 = V_N := 0$.

% Moreover, for any $j_t^*$ that achieves the minimum in \eqref{eq:Vt_def}, it holds that $\underline{k_t^*} < j_t^* \leq \overline{k_{t+1}^*}$, \sout{which implies that $P_t^*$ is stochastically smaller than $P_{t+1}^*$}. {\color{red} I'm confident that $\underline{k_t^*} \leq \underline{k_{t+1}^*}$ also holds, so that $\underline{k_t^*}$ doesn't have crossing issue.}
\end{theorem}
\begin{proof}
Denote the minimax objective function by 
\[O(\{P_n\}_{n=1}^N, P_y) = \sum_{n=1}^N \sum_{i=0}^{m-1} \sum_{b=1}^F w_{ib}^n \E_{p_n \sim P_n, Y \sim P_y}[S_{\tau_n,\theta_i}(p_n,Y) - S_{\tau_n,\theta_i}(f_b^n(t),Y)].
\]
In a similar vein to the proof of Lemma \ref{lem:minimax_sol_single_q}, we show $O(\{P_n^*\}_{n=1}^N, P_y^*) = \max_{P_y} O(\{P_n^*\}_{n=1}^N, P_y)$ then show $O(\{P_n^*\}_{n=1}^N, P_y^*) = \min_{P_n \in \Delta_m, n\in[N]} O(\{P_n\}_{n=1}^N, P_y^*)$.

Following the decomposition \eqref{eq:oppy_q}, $O(\{P_n\}_{n=1}^N, P_y)$ can be written as
\[
O(\{P_n\}_{n=1}^N, P_y) = \sum_{n=1}^N \sum_{i=0}^{m-1} \sum_{b=1}^F w_{ib}^n \E_{p_n \sim P_n, Y \sim P_y}[\one(Y \leq \theta_i < p_n) - \tau_n \one(\theta_i<p_n) - \one(Y \leq \theta_i < f_b^n(t))] + C_n,
\]
where $C_n =\tau_n \sum_{i=0}^{m-1} \sum_{b=1}^F w_{ib}^n \indi{\theta_i < f_b^n(t)}$ which does not depend on $P_n$ or $P_y$.

Using this decomposition, we prove $O(\{P_n^*\}_{n=1}^N, P_y^*) = \max_{P_y} O(\{P_n^*\}_{n=1}^N, P_y)$. Let $P_y(\frac{j}{m})=p_{yj}$ for $j=0,...,m$. Then, again using the notation from the proof of Lemma \ref{lem:jn_sort}, $O(\{P_n^*\}_{n=1}^N, P_y)$ is rephrased as the following:
\begin{align}\label{eq:O_Pstar_Py_decomp}
O(\{P_n^*\}_{n=1}^N, P_y) 
& = \sum_{n=1}^N \E_{p_n \sim P_n^*} \left[\sum_j p_{yj} \sum_{i=j}^{m-1} w_i^n \cdot \left( \indi{\theta_i < p_n} \one_F - \one_i^n \right)\right] + C(P_n^*),
\end{align}
where $C(P_n^*) = - \tau_n \E_{p_n \sim P_n^*} \sum_{i=0}^{m-1} w_i^n \cdot (\indi{\theta_i < p_n} \one_F - \one_i^n)$ only depends on $P_n^*$ and does not depend on $P_y$.

Plugging in $P_n^*$ and abbreviating $\indi{\theta_i < \hat{f}_i^n}$ with $\one_i^n$, the summand can be expressed as 
\begin{align*}
A_j^n := \E_{p_n \sim P_n^*} & \left[\sum_{i=j}^{m-1} w_i^n \cdot \left( \indi{\theta_i < p_n} \one_F - \one_i^n\right)\right] 
= \begin{cases}
-\sum_{i=0}^{j-1} w_i^n \cdot (\one_F-\one_i^n) - V_n + V_{n-1} & \quad j \leq k_n^* \\
- \sum_{i=j}^{m-1} w_i^n \cdot \one_i^n & \quad j > k_n^* \\
\end{cases} 
\end{align*}

This leads to 
\begin{align*}
A_j & := \sum_{n=1}^N A_j^n \\
&= \begin{cases}
-\sum_{n=1}^N \sum_{i=0}^{j-1} w_i^n \cdot (\one_F-\one_i^n) + V_0 & \quad j \leq k_1 \\
-\sum_{n=1}^\nu \sum_{i=j}^{m-1} w_i^n \cdot \one_i^n - \sum_{n=\nu+1}^{N} \sum_{i=0}^{j-1} w_i^n \cdot (\one_F-\one_i^n) + V_\nu & \quad  j \in (k_\nu, k_{\nu+1}], \,\, \nu \in [N-1] \\
-\sum_{n=1}^N \sum_{i=j}^{m-1} w_i^n \cdot \one_i^n + V_N & \quad j > k_N
\end{cases}
\end{align*}
Notice that $A_j$ is non-increasing on $j \leq k_1$ and non-decreasing on $j > k_N$. $A_0 = 0$ and $A_m = 0$ is the maximum value of $A_j$ in these ranges, respectively. Next, observe that $A_j \leq 0$ for all $j=0,1,...,m$ since $V_n$ is the minimum value that $\sum_{\nu=1}^n \sum_{i=j}^{m-1} w_i^n \cdot \one_i^\nu + \sum_{\nu=n+1}^{N} \sum_{i=0}^{j-1} w_i^{\nu} \cdot (\one_F-\one_i^{\nu})$ can achieve by its definition. It suffices to show that $A_{j_n^*} = 0$ for $n \in [N-1]$ since $P_y^*$ only puts mass on these points. Note that $A_{j_0} = A_0 = 0$ and $A_{j_N} = A_m = 0$.

From the first and second condition in the statement of Lemma \ref{lem:Pn_def}, we have $k_\nu < j_\nu^* \leq k_{\nu+1}$ if $j_{\nu-1}^* \neq j_\nu^*$ and $j_{\nu-1}^* = j_\nu^*= k_\nu$ if $j_{\nu-1}^* = j_\nu^*$, for $\nu \in [N-1]$. In the former case, we obtain $A_{j_\nu^*} = 0$ from $k_\nu < j \leq k_{\nu+1}$ case in the definition of $A_j$ above. In the latter case, $j_\nu^* = k_\nu$ does not fall into the range $(k_\nu, k_{\nu+1}]$. Still, from $k_{\nu-1} < j \leq k_{\nu}$ case ($j \leq k_1$ case if $\nu=1$) in the definition of $A_j$ , we have
\begin{align*}
A_{j_n^*} 
& = -\sum_{n=1}^{\nu-1} \sum_{i=j}^{m-1} w_i^n \cdot \one_i^n - \sum_{n=\nu}^{N} \sum_{i=0}^{j-1} w_i^n \cdot (\one_F - \one_i^n) + V_{\nu-1} \\
& =  - V_{\nu-1}(j_n^*) + V_{\nu-1} = -V_{\nu-1}(j_{n-1}^*) + V_{\nu-1} = 0,
\end{align*}
where $V_n(j)$ is defined as equation \eqref{eq:Vn_j_def}. Therefore, $A_{j_n^*} = 0$ for all $n \in [N-1]$ and $O(\{P_n^*\}_{n=1}^N, P_y^*) = \max_{P_y} O(\{P_n^*\}_{n=1}^N, P_y)$.

Next, we show that $O(\{P_n^*\}_{n=1}^N, P_y^*) = \min_{P_n \in \Delta_m, n\in[N]} O(\{P_n\}_{n=1}^N, P_y^*)$. This follows directly from the consistency of $S_{\tau,\theta_i}$ losses and the third condition in Lemma \ref{lem:Pn_def}. Any point in the support of $P_n^*$ is $\tau_n$-level quantile of $P_y^*$ since $P_y(Y \leq j_{n-1}^*/m) = \tau_n$ and $P_y(Y \geq j_n^*/m) = \tau_n$.

Lastly, we compute the value of the minimax program. Decomposition \eqref{eq:O_Pstar_Py_decomp} can be rewritten as 
\[
O(\{P_n^*\}_{n=1}^N, P_y) = \sum_{j} p_{yj} A_j + \sum_{n=1}^N C(P_n^*).
\]
Since $P_Y^*$ only puts mass on points $j/m$ such that $A_j = 0$, plugging in $P_y^*$ into the equation above yields
\begin{align*}
O(\{P_n^*\}_{n=1}^N, P_y^*) 
&= \sum_{n=1}^N C(P_n^*) = -\sum_{n=1}^N \tau_n \E_{p_n \sim P_n^*} \sum_{i=0}^{m-1} w_i^n \cdot (\indi{\theta_i < p_n}\one_F - \one_i^n) \\
&  = -\sum_{n=1}^N \tau_n (V_{n-1} - V_n) = -\sum_{n=1}^{N-1}(\tau_{n+1}-\tau_n)V_n \leq 0,
\end{align*}
where the second from the last equality used the fourth condition in Lemma \ref{lem:Pn_def}.
\end{proof}


\subsection{Sample efficient algorithm for omniprediction of multiple quantile level forecasts}\label{subsec:omni_multi_q_online}

  


\begin{theorem}
  \label{thm:omni-twoplayer_online_multi_q}
  For a set of quantile levels $\tau_1 < \tau_2 < ... < \tau_N$, let $\cs$ denote the class of scoring functions with representation \eqref{eq:multi_q_consistent_loss} with positive measures $\{H_{\tau_n}\}_{n=1}^N$ such that $H_{\tau_n}(\R) \leq 1, \forall n \in [N]$ (the set of \textit{all} scoring functions that are consistent for the multiple quantile levels under mild regularity conditions). Consider a class of competitors $\cf$ where each forecaster $f \in \cf$ provides $\{\tau_n\}_{n=1}^N$-level quantile forecasts $\{f^n(t)\}_{n=1}^N$. With the target values $Y_t$, forecaster predictions $f^n(t)$ for $t \in [T], n \in [N], f \in \cf$ each taking values in a discretized grid of size $m$, and the learning rate $\eta_w = O(\sqrt{\frac{\log (mN)}{T}})$ and $\eta_f = O(\sqrt{\frac{\log |\cf|}{T}})$, the output $\hat{P}_t = \{\hat{P}_t^n\}_{n=1}^N$ from Algorithm \ref{alg:omni-twoplayer_online_multi_q} achieves the following omniprediction error:
  \begin{align*}
  \sup_{S \in \cs,\, f \in \cf} \frac{1}{T} \sum_{n=1}^N \sum_{t=1}^T \E_{p \sim \hat{P}_t^n} [S(p, Y_t)] - \E[S(f^n(t), Y_t)]
  \leq O\left(\sqrt{\frac{\log mN|\cf|}{T}}\right)
  \end{align*}
\end{theorem}

\begin{proof}
  Straightforward extension of Theorem \ref{thm:omni-twoplayer_online_q}, using theorem \ref{thm:minimax_negative_multi_q}.
\end{proof}












\subsection{Separate omniprediction for each quantile level}\label{subsec:separate_multi_q}
\cite{fakoor2023flexible} showed that post-sorting and isotonic projection decrease the loss for \textbf{unweighted} weighted sums of pinball losses at different quantile levels. Unfortunately, post-sorting does not achieve the same for \textbf{weighted} sums of pinball losses.

Define the pinball loss at quantile level $\tau$ as $\rho_\tau(x) := \tau x\indi{x \geq 0} +(1-\tau)|x|\indi{x < 0} = (\tau - \indi{x < 0})x$. Then define the loss function 
\[
\ell(x_1, x_2) := \E_Y [w_\alpha \rho_\alpha (y-x_1) + w_\beta \rho_\beta (y-x_2)].
\] 
Consider the case $\alpha < \beta$. We would like $\ell(x_1, x_2) - \ell(x_2, x_1) \geq 0$ for $x_1 > x_2$. 
\paragraph{Counterexample} Let $Y \sim \text{Ber}(1/2)$ and $0 < x_2 < x_1 < 1$. Then, 
\begin{align*}
\ell(x_1, x_2) - \ell(x_2, x_1) 
& = w_\alpha \E_Y \left[(\alpha - \indi{y < x_1})(y-x_1) - (\alpha - \indi{y < x_2})(y-x_2)\right] + \\ 
& \qquad w_\beta \E_Y \left[(\beta - \indi{y < x_2})(y-x_2) - (\beta - \indi{y < x_1})(y-x_1) \right] \\
& = w_\alpha (\alpha - 1/2)(x_2 - x_1) + w_\beta (\beta - 1/2)(x_1 - x_2) \\
& = (x_1 - x_2)(w_\beta (\beta - 1/2) - w_\alpha (\alpha - 1/2))
\end{align*}
Note that the second equality uses the fact that $\indi{y < x_1} = \indi{y < x_2}$. When $\alpha$ and $\beta$ lie on the same side of $1/2$, there exist $(w_\alpha, w_\beta)$ for which the improvement in loss becomes negative.
Conversely, if $(w_\alpha, w_\beta)$ does not satisfy any of the above conditions, then there exists $(Y, x_1, x_2)$ such that post-sorting increases the loss.

\smallskip
If we solve the minimax program \textbf{separately for each quantile level}, quantile crossing may occur (or maybe NOT due to the structure of $f(X)$ that it is a valid probability forecast and thus no quantile crossing occurs). Unfortuately, the following proposition shows that post-sorting does not always reduce loss for \textbf{unequally} weighted sums of pinball losses.

\subsubsection{Characterizaion of weighted quantile loss that post-sorting does decrease the loss}
\begin{prop}
As above, define a weighted quantile loss as
\[
\ell(x_1, x_2) := \E_Y [w_\alpha \rho_\alpha (y-x_1) + w_\beta \rho_\beta (y-x_2)].
\]
If $(w_\alpha, w_\beta)$ satisfies one of the following conditions, then post-sorting does not increase the loss:
\begin{enumerate}
  \item $w_\alpha = w_\beta$
  \item $w_\alpha < w_\beta$ and $(1-\alpha)w_\alpha \geq (1-\beta)w_\beta$
  \item $w_\alpha > w_\beta$ and $\beta w_\beta \geq \alpha w_\alpha$.
\end{enumerate}
Conversely, if $(w_\alpha, w_\beta)$ does not satisfy any of the above conditions, then there exists $(Y, x_1, x_2)$ such that post-sorting increases the loss.
\end{prop}

\begin{proof}
\begin{align*}
\ell(x_1, x_2) - \ell(x_2, x_1) 
& = w_\alpha \E_Y \left[(\alpha - \indi{y < x_1})(y-x_1) - (\alpha - \indi{y < x_2})(y-x_2) \right] + \\ 
& \qquad w_\beta \E_Y \left[(\beta - \indi{y < x_2})(y-x_2) - (\beta - \indi{y < x_1})(y-x_1) \right] \\
& = w_\alpha E_Y \left[\left\{-\alpha(x_1 - x_2) + \indi{y < x_1}x_1 - \indi{y < x_2}x_2 - \indi{x_2 \leq y < x_1}y \right\}\right] + \\
& \qquad w_\beta E_Y \left[ \left\{\beta(x_1 - x_2) + \indi{y < x_2}x_2 - \indi(y < x_1)x_1 + \indi{x_2 \leq y < x_1}y \right\} \right] \\
& = \E_Y \left[\left\{-\alpha(x_1 - x_2) + \indi{x_2 \leq y < x_1}(x_1 - y) + \indi{y < x_2}(x_1 - x_2)\right\} \right] w_\alpha + \\
& \qquad \E_Y \left[\left\{\beta(x_1 - x_2) - \indi{x_1 \leq y < x_2}(x_1 - y) - \indi{y < x_2}(x_1 - x_2)\right\} \right] w_\beta 
\end{align*}
Denote the coefficient of $w_\alpha, w_\beta$ as $c_\alpha, c_\beta$ respectively. Notice that $c_\alpha + c_\beta = (\beta-\alpha)(x_1 - x_2) > 0$. Therefore $\ell(x_1, x_2) - \ell(x_2, x_1) > 0$ if $w_\alpha = w_\beta$ (Condition I).
Denote $A = \P(y < x_2)$, $B = \P(x_2 \geq y < x_1)$ and observe that $\E_Y \indi{x_2 \geq y < x_1}(x_1 - y) \in [0, B(x_1 - x_2)]$.
We now consider Condition II: $w_\alpha < w_\beta$ and $(1-\alpha)w_\alpha \geq (1-\beta)w_\beta$. 
\begin{align*}
\ell(x_1, x_2) - \ell(x_2, x_1) 
& = \left(w_\alpha(A-\alpha) + w_\beta(\beta-A)\right)(x_1 - x_2) - (w_\beta - w_\alpha)\E_Y \indi{x_2 \leq y < x_1}(x_1 - y) \\
& \geq -(w_\beta - w_\alpha)A + \beta w_\beta - \alpha w_\alpha - (w_\beta - w_\alpha)B(x_1 - x_2) \\
& \geq -(w_\beta - w_\alpha)A + (w_\beta - w_\alpha) - (w_\beta - w_\alpha)B(x_1 - x_2) \\
& \geq (w_\beta - w_\alpha)(1-A-B)(x_1 - x_2) \geq 0.
\end{align*}
Note that $\beta w_\beta - \alpha w_\alpha \geq w_\beta - w_\alpha$ follows from the condition and $A+B = \P(Y < x_1) \leq 1$.
Condition III is similar.
{\color{red} If we put this propoisition in the paper, add constructing counterexample for $(w_\alpha, w_\beta)$ that does not satisfy any of the above conditions.}
\end{proof}

Alternatively, we prove another class of loss functions where post-sorting does not increase the loss. Due to \cite{fissler2016multiq}, a loss function that is consistent for multiple quantile levels can be written as a weighted sum of consistent (proper) loss for each quantile levels, under mild regularity conditions:
\begin{align} \label{eq:multi_q_consistent_loss}
\ell(\{x_i\}_{i=1}^N, y) = \E_Y \sum_{i=1}^{N} w^i (\indi{y < x_i} - \tau_i)(g_i(x_i) - g_i(y)),  
\end{align}
where $g_i$ is a non-decreasing function.
\begin{prop}
If a loss function $\ell$ defined as \eqref{eq:multi_q_consistent_loss} satisfies $w_i = w$ and $g_i(x) = g(x)$ for all $i \in [N]$, then post-sorting does not increase the loss.
\end{prop}
\begin{proof}
Due to additivity of the loss function, it suffices to show $\ell(x_1, x_2, y) - \ell(x_2, x_1, y) \geq 0$ for $x_1 > x_2$ where $x_i$ corresponds to the $\tau_i$-level quantile estimate. Without loss of generality, assume $w=1$.
\begin{align*}
\ell(x_1, x_2, y) & - \ell(x_2, x_1, y) 
= \E_Y [ (\indi{y < x_1} - \tau_1)(g_1(x_1) - g_1(y)) - (\indi{y < x_2} - \tau_1)(g_1(x_2) - g_1(y)) + \\
& \qquad (\indi{y < x_2} - \tau_2)(g_2(x_2) - g_2(y)) - (\indi{y < x_1} - \tau_2)(g_2(x_1) - g_2(y))] \\
& = \E_Y [\tau_1(g_1(x_2)-g_1(x_1)) + \indi{y < x_1}g_1(x_1) - \indi{y < x_2}g_1(x_2) - \indi{x_2 \leq y < x_1}g_1(y) + \\
& \qquad \tau_2(g_2(x_1)-g_2(x_2)) + \indi{y < x_2}g_2(x_2) - \indi{y < x_1}g_2(x_1) + \indi{x_2 \leq y < x_1}g_2(y)]
\end{align*}
Now using the fact that $g_1 \equiv g_2 \equiv g$, most of the terms cancel out and we have 
\begin{align*}
\ell(x_1, x_2, y) - \ell(x_2, x_1, y) = (\tau_2 - \tau_1)(g(x_1) - g(x_2)) \geq 0.
\end{align*}
{\color{blue} Just to add, in general case with $w_i=w$ but $g_i$ can be different, the difference of loss is as following:
\begin{align*}
\ell(x_1, x_2, y) - \ell(x_2, x_1, y) 
& = \tau_2(g_2(x_1)-g_2(x_2)) - \tau_1(g_1(x_1)-g_1(x_2)) - \int_{x_2}^{x_1} (\P(Y < x))(g_2'(x) - g_1'(x)) dx \\
& = (\tau_2 - \tau_1)(g_1(x_1) - g_1(x_2)) - \int_{x_2}^{x_1} (\P(Y < x) - \beta)(g_2'(x) - g_1'(x)) dx  
\end{align*}
Note that $g_i(x_1) - g_i(x_2) \geq 0$ since $g_i$ is non-decreasing.
}
\end{proof}



